\chapter{Research Directions}
\label{ch:research-directions}

Each item below is attributed to where it came from. Items marked \textbf{flagged} were named as
open by the speaker. Items marked \textbf{inferred} are editorial, drawn from placing two lectures
next to each other, and should be treated as untested conjecture rather than as anything the
speakers endorsed.

\section*{Problems the speakers explicitly left open}

\paragraph{Measuring what we care about when building a representation. Flagged, Pascanu
(Chapter~\ref{ch:07-to-be-figured-out}).} He argues that non-i.i.d.\ structure cannot be exploited
because we do not know how to measure representation quality during training, and separately because
learning lacks compositionality and locality. The measurement problem is the tractable half: it is a
concrete request for a training-time signal that predicts downstream representation quality.

\paragraph{Evaluating continual learning at all. Flagged, Lyle
(Chapter~\ref{ch:06-continual-learning}).} She states the meta-problem plainly and neither option is
acceptable: a pre-specified task sequence is artificial, a naturally continual application is
uncontrolled. Any result in the field inherits this weakness, so a defensible evaluation protocol
would be a contribution in itself.

\paragraph{In-context learning as a continual learning mechanism. Flagged, Lyle.} Her closing slide
calls this an exciting paradigm that brings its own challenges, without saying which. Given the rest
of her lecture, the obvious question is whether in-context adaptation escapes the plasticity problem
or merely relocates it into the context window.

\paragraph{What else admits a learned estimator. Flagged, Cho
(Chapter~\ref{ch:05-statistics-estimation}).} Asked directly on a slide. He has done population
standard deviation, mutual information, Bayesian predictive distributions and causal effects. The
programme's requirement is a simulator of problems with known answers, which is the real filter on
what comes next.

\paragraph{Whether a world model aligns the way vision and language models do. Flagged, Isola
(Chapter~\ref{ch:02-representational-geometry}).} Posed as an open question on a slide. His
convergence results are for vision and language; whether a model trained on interaction and dynamics
lands on the same kernel is untested and would discriminate between the three hypotheses he offers
about what "same" means.

\paragraph{Generalisation in diffusion models, and flow maps. Flagged, Dieleman
(Chapter~\ref{ch:03-diffusion-models}).} His closing slide lists what he did not cover:
generalisation versus memorisation, integrating diffusion models via flow
maps~\citep{dieleman2026flowmaps}, reward-based steering and finetuning, and multimodal generative
models. The flow-maps thread is also flagged in the personal notes.

\paragraph{Quantized training, and extending the linearity theorem. Flagged, Alistarh
(Chapter~\ref{ch:10-efficient-learning}).} He names quantized training as the frontier and says big
open questions sit at the intersection of representations, optimisation and compute. There is also a
precise gap on the slide: the linearity theorem holds only for methods that do not update weights, so
it covers AWQ but not GPTQ, which is awkward because GPTQ's compensating update is its whole
mechanism.

\paragraph{Building faithful evaluators, and learning a Galois action. Flagged,
Tapu\v{s}kovi\'{c} (Chapter~\ref{ch:08-ml-for-mathematics}).} His stated conclusion is that search is
not the bottleneck and constructing a cheap, exact, ungameable score is the research. He also asks
directly whether the action of the Cosmic Galois Group on Feynman Mellin motives can be learned, and
observes that machine certification stops exactly where the deep objects begin, which is a diagnosis
of where formalisation effort should go.

\paragraph{Gauge equivariance without parallel transport. Flagged, Veli\v{c}kovi\'{c}
(Chapter~\ref{ch:11-geometric-deep-learning}).} Conditioning across arbitrarily related pairs of
gauges is described as very hard, and the working solution sidesteps it by transporting features
first. Whether the harder formulation buys anything is open.

\paragraph{Why real-world reinforcement learning fails. Flagged, Deac
(Chapter~\ref{ch:04-intro-reinforcement-learning}).} Her own live bandit collected zero votes over 53
rounds, and the closing slide enumerates the causes: reward confounded with content,
non-stationarity, sparsity, noise, varying turnout, delayed and batched feedback, and a false
independence assumption across action dimensions. That list is a better research agenda than most
papers' future-work sections.

\section*{Directions inferred here, not reported}

\paragraph{Give causally-masked transformers a spectrum. Inferred, from Stankovi\'{c} and
Veli\v{c}kovi\'{c}.} Stankovi\'{c} notes that a causal attention mask is strictly upper triangular and
therefore nilpotent, and Veli\v{c}kovi\'{c} argues a transformer \emph{is} a GNN on that graph. Putting
the two together, every decoder-only transformer is a graph whose adjacency spectrum has collapsed,
and her zero-padding construction is a candidate for restoring it: add a return path with enough zero
tokens that the induced feedback cannot reach real tokens within the depth of the network. Whether
the restored spectrum says anything useful about attention is entirely untested, and the construction
was designed for graph filters rather than learned attention, so this may not survive contact.

\paragraph{An experiment that would separate Lyle from Pascanu. Inferred.} Lyle claims plasticity
loss is a fixable optimiser and architecture fault; Pascanu claims gradient-based credit assignment
structurally needs i.i.d.\ data. Apply Lyle's full remedy, normalisation, norm control, avoiding
ill-conditioned targets, and then test not whether training remains stable, which she has shown, but
whether knowledge \emph{composes}: does a model trained on tasks $A$ then $B$ solve compositions of
$A$ and $B$ better than one trained on $B$ alone? Lyle's account predicts that stability suffices;
Pascanu's predicts that composition still fails. No such experiment appears in either lecture.

\paragraph{A learned representational similarity metric. Inferred, from Cho and Isola.} Isola says
unlearned distance functions have no future; Cho shows how to learn an estimator whenever problems
with known answers can be simulated. Representational similarity is a candidate: train an estimator
on pairs of representations related by known transformations, with the transformation as the label.
The hazard is the circularity flagged in Chapter~\ref{ch:themes}: convergence results measured with
a fitted metric are hard to interpret, and confronting that directly might be the more valuable
output.

\paragraph{Model the compensating update in the linearity theorem. Inferred, from Alistarh.} The
theorem's disclaimer excludes weight-updating methods. GPTQ's update is not arbitrary though: it is
the closed-form $\dw^{*}$ from the inverse Hessian. That suggests the theorem might extend by
carrying the update through the per-layer error term rather than by treating it as unmodelled noise.
Speculative, and whether the assumptions survive has not been checked.

\begin{aside}
Every item in the second list is a guess made by putting two lectures side by side, and none of them
was checked against the literature. Anyone who finds that one of them has already been done, or is
already known not to work, is better informed than these notes are.
\end{aside}
